\chapter{Background: Bottlenecks in BSP Systems}

% Maybe note ML differences as a separate section

We now discuss key concepts and challenges in the bulk-synchronous parallel (BSP) paradigm and
introduce state-of-the-art approaches for different things.
While our initial work focused on scientific applications, a review of literature reveals similar
challenges in distributed machine learning training.
These common challenges inform our later work, and this section discusses background and
outstanding challenges in both domains.

\section{The High-Performance Computing Environment}

\paragraph{System Architecture}
Large-scale HPC clusters typically consist of O(10,000) nodes~\cite{:TOP500ListJune} connected by a
high-performance fabric~\cite{Pfist2001:Infiniband,Rowet2022:Slingshot,Metz2024:UltraEthernet}.
Each node typically consists of CPUs and typically multiple GPUs, which are increasingly important
for scientific computing~\cite{DeBar2014:HPCGPU,Snell2017:HPCGPU}.
Storage is traditionally provided by parallel filesystems~\cite{Braam2002:Lustre} but object
stores~\cite{Henne2020:DAOS} and richer commercial offerings~\cite{2024:Hammerspace,:VASTPlatform}
are increasingly common~\cite{Glenn2025:NoPFS}.

\paragraph{Applications}
Scientific applications target fundamental questions across domains such as cosmology and climate
science, problems that require simulating complex physical phenomena over extended spatial and
temporal scales.
These applications can occupy the majority of a supercomputer's resources for months at a
time~\cite{Das2023:LargeScale, Stock2024:LargeScale,Gary2023:LANLPlatformPlanning}.
From a systems perspective, they can be simplified into two common archetypes: \emph{particle
  codes}~\cite{Bird2022:VPIC2.0} and \emph{mesh codes}~\cite{Dubey2014:Survey}.
Particle codes track the state and interactions of billions individual discrete entities as they
move through the simulation domain.
Mesh codes discretize the spatial domain into grids and compute the evolution of properties of
interest at each grid point.

% todo: ML training?

\paragraph{The BSP Model}
The BSP model consists of alternating compute and synchronization phases.
Execution begins with decomposition of the problem into independent subproblems assigned to
workers, and proceeds as discrete timesteps.
Each timestep consists of a compute phase followed by a global synchronization phase, where
synchronous collectives ensure all ranks advance together—a requirement for correctness.
Every few timesteps, a checkpoint is triggered for fault tolerance offline analysis.
Checkpointing frequency increases with application scale as intra-checkpoint failure probability
rises, creating periodic I/O-bound operations that stress the storage system.

\paragraph{Distributed ML Training}
Distributed ML training represents a similar and increasingly important cluster workload, occupying
tens of thousands of GPUs and running for months~\cite{:LlamaScale,Rahma2024:MLScale,Gratt2024:Llama3Scale}.
The cluster architectures are comparable, though ML clusters often add a \emph{frontend} Ethernet
network for general datacenter connectivity~\cite{Gangi2024:RDMA}.
While the parallelization strategies differ---ML employs a mix of data, tensor, and pipeline
parallelism---the execution model still relies on the same core BSP primitives of synchronous
collectives~\cite{Gangi2024:RDMA} and periodic checkpointing~\cite{Strat:CheckpointCheckpointsLarge}.


\section{Data Management}


Ingestion efficiency essential.
WAF of 1X ideal and essential to ingest at max storage system bandwidth.
Databases are not ideal.

Special-purpose systems used for in-situ aggregation and visualization.
In-situ indexing and reordering less widespread but demonstrated to be viable.
Deltafs (insitu hash partitioning).
But too much reliance on bitmap indices etc, as partitioning is expensive.
Partitioning for range queries is hard: this currently uses post-processing sorts, which need to
read all data again and write them, requiring extra hardware and delaying time to insight.

The biggest missing piece in the in-situ toolking is range partitioning.

\section{Compute Efficiency}

\subsection{Workload Characterization}
Reliance on aggregate metrics.
Did this improve runtime overall?

Fine-grained effects hard to untangle, but esesential, as our work shows.
Only available with tracing.

Trace analysis model: a small snapshot of execution results in a large trace.
Published and analyzed by other groups.
This makes learnings slow and disconnected from hardware.
Those who analyze traces are not able to distinguish fundamental characteristics from tunable
runtime artifacts.

\subsection{Perf Tools}

Perf tools can be sampling, profiling, or tracing.
Former emit summaries.
Latter emits lots of raw data in formats like OTF2.
These are not analytical formats and are not amenable to analytics.

ML: Pytorch/Kineto traces.
Clusters starting to eclipse HPC.
Similar runtime challenges.

Solutions for perf challenges are hardcoded assertions in different tooling.








